\documentclass{article}
\usepackage[margin=28mm]{geometry}
\usepackage{amsmath,amssymb}
\usepackage{xurl}
\usepackage[hidelinks]{hyperref}
% Shared commands for notes in docs/paper-gaps/.

\DeclareUnicodeCharacter{2080}{\ifmmode{}_{0}\else\textsubscript{0}\fi}
\DeclareUnicodeCharacter{2081}{\ifmmode{}_{1}\else\textsubscript{1}\fi}
\DeclareUnicodeCharacter{2082}{\ifmmode{}_{2}\else\textsubscript{2}\fi}
\DeclareUnicodeCharacter{2099}{\ifmmode{}_{n}\else\textsubscript{n}\fi}

\newcommand{\Lean}{\textsc{Lean}}
\ExplSyntaxOn
\NewDocumentCommand{\leanid}{m}{
  \ifmmode
    \textsf{\detokenize{#1}}
  \else
    \group_begin:
    \urlstyle{sf}
    \tl_set:Nn \l_tmpa_tl {#1}
    \tl_replace_all:Nnn \l_tmpa_tl {\_} {_}
    \exp_args:NV \path \l_tmpa_tl
    \group_end:
  \fi
}
\ExplSyntaxOff
\providecommand{\tr}{\operatorname{tr}}
\newcommand{\tnleanIssueBase}{https://github.com/LionSR/TNLean/issues/}
\newcommand{\tnleanPrBase}{https://github.com/LionSR/TNLean/pull/}
\newcommand{\tnleanDocBase}{https://sirui-lu.com/TNLean/docs/}
\newcommand{\ghissue}[1]{\href{\tnleanIssueBase#1}{GitHub issue \##1}}
\newcommand{\ghpr}[1]{\href{\tnleanPrBase#1}{PR \##1}}
\newcommand{\doclink}[1]{\href{\tnleanDocBase#1.html}{\path{#1}}}

% Dirac notation and the identity operator, as in blueprint/src/macros/common.tex.
% \providecommand keeps note-local definitions authoritative.
\usepackage{dsfont}
\providecommand{\ket}[1]{|#1\rangle}
\providecommand{\bra}[1]{\langle #1|}
\providecommand{\braket}[2]{\langle #1 | #2 \rangle}
\providecommand{\ketbra}[2]{|#1\rangle\!\langle #2|}
\providecommand{\Id}{\mathds{1}}
\providecommand{\one}{\mathds{1}}
\providecommand{\C}{\mathbb{C}}
\providecommand{\MN}[1]{M_{#1}(\C)}
\providecommand{\MD}{M_D(\C)}

% Verdict marker: \gapnote{<kind>}{<status>}. Kind is the policy
% classification (clarification, local-correction, scope-restriction,
% unfaithful, false-source, open-gap); status is open, wip (elimination
% underway), resolved, or historical. Machine-read by the site index;
% prints a verdict line.
\newcommand{\gapnote}[2]{\par\noindent\textbf{Verdict:} #1 (#2).\par}

\providecommand{\gapnote}[2]{\par\noindent\textbf{Verdict:} #1 (#2).\par}
\providecommand{\leanid}[1]{\textsf{\detokenize{#1}}}
\title{Positivity of the Matrix Product Form at Every Temperature}
\author{}
\date{}
\begin{document}
\maketitle
\gapnote{local-correction}{resolved}

\section{The source statement}
Hastings, \emph{Solving gapped Hamiltonians locally}~\cite{Hastings2006Solving},
\href{https://arxiv.org/abs/cond-mat/0508554v3}{arXiv:cond-mat/0508554v3},
source \path{proj3new.tex}, lines 1303--1314, sets
$\rho(\beta,l_{\mathrm{proj}})=\rho(\beta/n,l_{\mathrm{proj}})^n$ for
$\beta>\beta_0$ (eq.~(mplt)), with $n$ the smallest even integer such that
$\beta/n\le\beta_0$, and states: ``By picking $n$ even we guarantee that
$\rho(\beta,l_{\mathrm{proj}})$ is positive definite.''

\section{The point at issue}
The truncated high-temperature form $A=\rho(\beta/n,l_{\mathrm{proj}})$ is
Hermitian. For even $n=2k$,
\begin{align}
  A^n=(A^k)^\dagger A^k,
  \label{eq:hastings05_mplt_positivity_square}
\end{align}
so $\langle v,A^nv\rangle=\lVert A^kv\rVert^2\ge0$ for every vector $v$: the
power is positive semidefinite. It is positive definite exactly when $A^k$ has
trivial kernel, that is, when $A$ is invertible. The source does not show that
the truncated form $A$ is invertible. It approximates $e^{-\beta H/n}$, which
is invertible, only in the trace-norm sense of lines 1281--1299, and an
approximation in trace norm need not be invertible.

\section{The statement proved}
The formalization proves that the even power is positive semidefinite, which is
what the evenness of $n$ guarantees. Positive definiteness holds under the
additional hypothesis that $\rho(\beta/n,l_{\mathrm{proj}})$ is invertible,
which is not established. Nothing later in the source's argument uses strict
positivity: the purpose of the even power is that
$\rho(\beta,l_{\mathrm{proj}})$ be a positive operator, as a density matrix
should be, and semidefiniteness suffices for that.

\section{Formalization}
The positive semidefiniteness of even powers is
\leanid{TNLean.Lattice.ClusterExpansion.posSemidef_pow_of_even}, and together
with the matrix product form it is
\leanid{TNLean.Lattice.HasProductForm.pow_of_even}.
\bibliographystyle{unsrt}
\bibliography{references}
\end{document}
